\documentclass[a4paper,11pt]{jsarticle}
\usepackage{amsmath}
\usepackage{booktabs}
\title{TeX Sync View デモ文書}
\author{}
\date{}
\begin{document}
\maketitle

\section{はじめに}
この文書は、拡張機能 TeX Sync View の動作を紹介するためのダミーです。
左のエディタで文章を選択すると、右の PDF の対応する位置が黄色でマークされます。
逆に PDF をダブルクリックすると、エディタの該当行へ移動します。

エディタと PDF のスクロールは双方向に連動します。
どちらを動かしても、もう一方が同じ箇所を追いかけます。

\section{数式}
円周率に関する有名な式を例にします。
\begin{equation}
  \frac{\pi}{4} = 1 - \frac{1}{3} + \frac{1}{5} - \frac{1}{7} + \cdots
\end{equation}
ガウス積分は次のとおりです。
\begin{equation}
  \int_{-\infty}^{\infty} e^{-x^2}\,dx = \sqrt{\pi}
\end{equation}
複数行にわたる式も、範囲をまとめて選択すれば全体がマークされます。
\begin{align}
  (a+b)^2 &= a^2 + 2ab + b^2 \\
  (a-b)^2 &= a^2 - 2ab + b^2
\end{align}

\section{表}
表の中の行も、対応する位置が正しくマークされます。
\begin{table}[h]
  \centering
  \begin{tabular}{lcc}
    \toprule
    項目 & 値 & 単位 \\
    \midrule
    長さ & 12.5 & cm \\
    重さ & 3.2 & kg \\
    温度 & 21.0 & ℃ \\
    \bottomrule
  \end{tabular}
  \caption{デモ用の表}
\end{table}

\section{おわりに}
選択範囲を広げるほど、マークも広がります。
自分の文書で試してみてください。
\end{document}
